For the two root tests, let \(g\) be the chosen level prime, let \(r\) be the rank, and put \(N=r\). For a degree-\(j\) label, the order of its multiplicative group and its \(\ell\)-primary part are
\[
  A_j=g^j-1,\qquad A_{j,\ell}=\ell^{v_\ell(g^j-1)}.
\]
 The
squarefree mass of labels that contain specified distinct
irreducibles of degrees $j_1,\ldots,j_u$ is at most
$\prod_i A_{j_i}^{-1}$.

Strip off the specified factors and sum
the remaining weights, which have total at most one in the
remaining degree. This last assertion also follows without the
character formula: if $I_j$ is the number of irreducibles of
degree $j$ with nonzero constant term, then $jI_j\le A_j$ by
counting their distinct roots in $\mathbb F_{g^j}^\times$.
Thus the mass of polynomials with factor counts $m_j$ is at
most $\prod_j(j^{m_j}m_j!)^{-1}$, the probability of that cycle
count under a uniform permutation. This proves domination for
any event specified by the factor degrees.

\begin{lemma}[Mass of the four excluded finite types]
\label{lem:finite-type-event-bounds}
Use the squarefree subprobability measure and the hypotheses of
\eqref{eq:row-bad-types}. Write \(m_{\rm pair},m_{\rm scal},m_{\rm div},m_{\rm cyc}\)
for the probabilities of a repeated prime-to-\(\ell\) root part, a match
with a prescribed scalar or its inverse, a common proper divisor degree,
and absence of the prescribed large prime degree, respectively. Then
\[
 \begin{aligned}
 m_{\rm pair}&\le 23f\log(g)/\sqrt g,&
 m_{\rm scal}&\le2fm\left(\frac1{4\min(q,d')}+\frac{1+\log r}{g-1}\right),\\
 m_{\rm div}&\le\tfrac r2(7/8)^f,&
 m_{\rm cyc}&\le\exp\!\left(-\frac f{5\log r}\right).
 \end{aligned}
\]
The pair test ranges over the nonexcluded \(\ell\), and the last two
bounds use the product over the \(f\) residue factors. The separate
summand \(3f/g\) in \eqref{eq:row-bad-types} allows for the missing
squarefree mass.
\end{lemma}

\begin{proof}
\proofstep{Repeated roots in the squarefree mass.}
When two different roots of a single irreducible of degree $j$
have the same prime-to-$\ell$ part, they are $z,z^{g^a}$ with
$1\le a<j$. Their equality requires $z^{g^a-1}$ to lie in the
$\ell$-primary group. In the cyclic group of order $A_j$, this
allows at most $(g^{\gcd(a,j)}-1)A_{j,\ell}$ choices of $z$.
For a nonexcluded prime, $A_{j,\ell}\le g^{j/4}$ by its
exclusion test.

Since \(\gcd(a,j)\le j/2\), the root count is bounded by
\[
  (g^{\gcd(a,j)}-1)A_{j,\ell}\le g^{j/2}g^{j/4}=g^{3j/4}.
\]
 An equality between distinct roots can occur only
for $\ell\mid A_j$. There are at most $j\log_2 g$ such primes.
Every represented irreducible has at least $j$ winning roots,
because the test is invariant under Frobenius. Dividing the root
count by $j$ and then by $A_j$ bounds the mass for one factor by
\[
 \frac{4096}{4095}\log_2 g
       \sum_{j\ge2}(j-1)u^j,
 \qquad u=g^{-1/4}\le1/8.
\]

For a pair of distinct irreducibles of degrees $i,j$ use the
cyclic groups on the two roots. The common prime-to-$\ell$
part lies in the group of order $g^{\gcd(i,j)}-1$, and the two
primary parts give at most $A_{i,\ell}A_{j,\ell}$ possibilities.
Thus the number of ordered root pairs is at most $g^{3(i+j)/4}$.
The union over the prime factors of $A_i A_j$, followed by
stripping the two factors from the label, bounds the mass by
\[
 \left(\frac{4096}{4095}\right)^2\log_2 g
                \sum_{i,j\ge1}(i+j)u^{i+j}.
\]
The sums are $u^2/(1-u)^2$ and $2u^2/(1-u)^3$. Together they
fit $23\log g/\sqrt g$, with a constant independent of $N$.
This also bounds the whole-place failure by the union over places.
\label{cor:paired-type-mass-1}

\proofstep{A fixed scalar and its root count.}
For a fixed scalar in $\mathbb F_g^\times$, the number of
roots in $\mathbb F_{g^j}^\times$ with matching prime-to-$\ell$
part is at most $A_{j,\ell}$. It suffices to take the union over
primes dividing $A_j$. If a match occurs at another prime the
root is the scalar itself, which is already counted. Divide by
$j$ for the number of roots of each irreducible, and by $A_j$
for the mass of labels containing it.
For fixed \(j,\ell\), these two divisions bound the mass from a fixed scalar by
\[
  \frac{1}{A_j}\,\frac{A_{j,\ell}}{j} .
\]
 At $j=1$, the sum of the
primary parts of $g-1$, divided by $g-1$, is at most
$3/(16\min(q,d'))$. Indeed the three parts for $2,q,d'$ have
size at least $16,q,d'$, and the product of any two is at least
$16\min(q,d')$.

For previous sum \(S\ge2\) and product \(P\), adjoining a primary part
\(c\ge2\) gives
\[
 \frac{S+c}{cP}\le\frac S P,
 \qquad cS-(S+c)=(c-1)(S-1)-1\ge0 .
\] At $j\ge2$, the product of the parts of
$A_j$ supported on $g-1$ is at most $j(g-1)$ by the lifting
formula for powers, also at two since $g\equiv1\pmod4$.
The product of the other parts is at most $A_j/(g-1)$.
In each group the sum of its primary parts is at most the
product. Thus the degree $j$ contribution is at most
$(g-1)/A_j+1/(j(g-1))$. Summing the resulting estimate uses both
\[
  \frac{g-1}{g^j-1}=\frac{1}{1+g+\cdots+g^{j-1}}\le g^{1-j},\qquad \sum_{j=2}^N\frac1j\le\log N.
\]
 This proves the scalar term
in \eqref{eq:row-bad-types} after summing over the $2fm$
scalars and inverse scalars.
\label{cor:paired-type-mass-2}

\proofstep{The divisor-degree and cycle tests.}
A divisor of degree $a$ for a squarefree label corresponds to
an invariant $a$-subset of a permutation in the degree-count
bound. For $0<a<N$, count stable $a$-subsets with weight eight,
except give weight four when the permutation fixes a point on
each side. For a fixed subset \(A\), let \(E_A\) be the event that
it is stable and \(F_A\) the event of a fixed point on
each side. Given \(E_A\), the two restrictions are
independent uniform permutations. The derangement formula
gives probability at least \(1/2\) of a fixed point on
each side, including a side of size one. Thus the mean
contribution from this subset is
\[
 \mathbb P(E_A)(8-4\mathbb P(F_A\mid E_A))\le7\mathbb P(E_A).
\]

When some subset is
stable the total weight is at least eight. In the weight-four
case, swap a fixed point on each side to obtain a different
stable subset of the same size. Summing the stabilization
probability over subsets gives one. Hence the chance of a
stable $a$-subset is at most $7/8$. Independence of the places
and a union over $1\le a\le N/2$ give the fourth term of
\eqref{eq:row-bad-types}.

A permutation can have at most one cycle of prime length in
$(N/2,N-3]$. For each such \(p\), the choice of its support, its cyclic ordering and the permutation on the complement gives
\[
  \Pr(\text{a cycle of length }p)=\binom Np\frac{(p-1)!(N-p)!}{N!}=\frac1p.
\]
 The sum of these reciprocals is at least
$1/(5\log N)$ for $N\ge100$. For $100\le N<16900$ use the sieve to count primes through
$16899$, and sum the lower bounds $10^{-12}\lfloor10^{12}/p\rfloor$.
Multiplying by $5\lfloor\log_2 N\rfloor\cdot0.6931471803$ gives
at least one for every such $N$, using only integer arithmetic.

For larger
$N$, write $n_1=\lfloor N/2\rfloor$, $s_1=\lfloor\sqrt N\rfloor$
and $t=\#\{p:N/2<p\le N-3\}$. The central binomial lower
bound $4^{n_1}/n_1$ and its prime factorization give
\[
 \frac{N-1}{6}\log 4 \le
                 \big(\pi(\sqrt{2n_1})+t+4\big)\log N.
\]
Indeed the primes at most $\sqrt{2n_1}$ cost at most $2n_1$
each, those between $\sqrt{2n_1}$ and $2n_1/3$ cost at most
the primorial up to $2n_1/3$, bounded by $4^{\lfloor2n_1/3\rfloor}$,
and those between $2n_1/3$ and $n_1$ do not divide the binomial.
At most three primes near the upper endpoint were omitted from
the list for $t$.

Now $\pi(\sqrt{2n_1})\le(s_1+9)/3$ by counting
classes coprime to $6$, and $\log N\le(s_1+1)/64+7.70407$ by
$\log u\le u-1$ applied to $u=(s_1+1)/128$. With $s_1\ge130$,
expanding the two sides gives
\[
 (s_1+21)((s_1+1)/64+7.70407)+0.693148\le0.093147s_1^2.
\]
Since $\log 2>0.69314718$, substitution gives $5t\log N\ge N$.
Each reciprocal in the sum is at least $1/N$. The cycle-count domination gives
$1-1/(5\log N)$ as an upper bound for the single-place mass
without one of the prescribed factors. Raising to the $f$th
power bounds the final event in \eqref{eq:row-bad-types}.
\label{cor:paired-type-mass-3}
\end{proof}
